\chapter*{复数}
\subsection*{复数的概念与代数运算}
\begin{enumerate}
    \item (2024$\cdot$崇明区一模$\cdot$3)若复数 $z = {m}^{2} - 4 + \left( {m + 2}\right) {i}$ ( ${i}$ 为虚数单位)是纯虚数，则实数 $m$ 的值为\blkda[2]{$2$}.
    \item (2025$\cdot$浦东新区一模$\cdot$3)已知复数 ${z}_{1} = 3 + {i},{z}_{2} = a + 4{i},a \in  \mathbf{R}$ ，若 ${z}_{1} - {z}_{2}$ 为纯虚数，则 $\left| {z}_{2}\right|  =$ \blkda[2]{$5$}.
    \item (2024 青浦区二模$\cdot$3)已知复数 $z = \dfrac{5{i}}{1 - {i}}$ ，则 ${Im}z =$ \blkda[2]{$\dfrac{5}{2}$}.
    \item (2025$\cdot$上海春季卷$\cdot$3)已知复数 $z = \dfrac{2 + {i}}{{i}}$ ，其中 ${i}$ 为虚数单位，则 $\left| z\right|  =$ \blkda[2]{$\sqrt{5}$}.
    \item (2024$\cdot$崇明区二模$\cdot$4)若复数 $z$ 满足 ${i}z = 1 + {i}$ ( ${i}$ 为虚数单位)，则 $z =$ \blkda[2]{$1-i$}.
\end{enumerate}
\subsection*{复数的共轭}
\begin{enumerate}
    \item (2024$\cdot$上海春季卷$\cdot$3)若复数 $z$ 满足 $\dfrac{z}{1 + {i}} = {i}$ ( ${i}$ 为虚数单位)，则 $\overline{z} =$ \blkda[2]{$- 1 - {i}$}.
    \item 已知集合$P = \{\omega | \omega = z + \overline{z}, z \in \mathbf{C}\}, Q = \{\omega | \omega = z - \overline{z}, z \in \mathbf{C}\}$，则$P \cap Q =$\blkda{$\{0\}$}
    \begin{jieda}
        $P$对应着实轴上的点，$Q$对应虚轴上的点，因此两者交集为$\{0\}$，对应着原点.
    \end{jieda}
    \item 已知下列四个命题：
    \begin{dingautolist}{172}
        \item 若两个虚数$z_1, z_2$的和是实数，则$z_1, z_2$是共轭虚数
        \item 若$z_1, z_2$是共轭复数，则$z_1, z_2$的差是纯虚数
        \item 若复数$z - \overline{z} = 0$，则$z$是实数
        \item 若复数$z + \overline{z} = 0$，则$z$是纯虚数
    \end{dingautolist}
    其中正确的命题有\blkda{\ding{174}}
    \begin{jieda}
        \begin{dingautolist}{172}
            \item 错误，和是实数说明$z_1, z_2$的虚部互为相反数，但实部未必相等
            \item 错误，若$z_1, z_2 \in \mathbf{R}$，则$z_1 = z_2$，因此$z_1 - z_2 = 0$
            \item 正确，$z - \overline{z} = 0$说明$z$的虚部为0，因此$z$实数
            \item 错误，反例：$z = 0$
        \end{dingautolist}
    \end{jieda}
    \item (2025$\cdot$普陀区一模$\cdot$3)设 ${i}$ 为虚数单位，若复数 $z$ 满足 $z \cdot  \overline{z} + z - \overline{z} = 4 + 2{i}$ ，则 $\left| z\right|  =$ \blkda[2]{$2$}.
    \item (2024$\cdot$上海秋季卷$\cdot$9)已知虚数 $z$ ，其实部为 1，且 $z + \dfrac{2}{z} = m\left( {m \in  \mathbf{R}}\right)$ ，则实数 $m$ 为\blkda[2]{$2$}.
    \begin{jieda}
        $z + \dfrac{2}{|z|^2} \cdot \overline{z} = m\left( {m \in  \mathbf{R}}\right)$，故必有$\dfrac{2}{|z|^2} = 1$，$m = 2Rez = 2$.
    \end{jieda}
\end{enumerate}

\subsection*{复数的模的几何意义}
\begin{enumerate}
    \item (2025$\cdot$杨浦区二模$\cdot$7)已知复数 $z$ 满足 $\left| {z - 1 + {i}}\right|  = 1$ ，其中 ${i}$ 为虚数单位，则 $\left| z\right|$ 的最小值为\blkda[2]{$\sqrt{2} - 1$}.
    \begin{jieda}
        在复平面内, $\left| {z - 1 + {i}}\right|  = 1$ 表示复数 $z$ 对应的点 $Z$ 与复数 $1 - {i}$ 对应点 $A\left( {1, - 1}\right)$ 的距离为 1,因此点 $Z$ 的轨迹是以 $A\left( {1, - 1}\right)$ 为圆心,1 为半径的圆, $\left| z\right|$ 表示点 $Z$ 到原点 $O$ 的距离，所以 $\left| z\right|$ 的最小值为 $\left| {OA}\right|  - 1 = \sqrt{2} - 1$.
    \end{jieda}
    \item 如果复数$z$满足$|z+i| + |z-i| = 2$，那么$|z+i+1|$的最小值是\blkda{1}
    \begin{jieda}
        根据复数的几何意义可知，$|z+i| + |z-i| = 2$即$z$在复平面上的点$Z$在$(0, 1), (0, -1)$两点间的线段上（含端点），$|z+i+1|$即$Z$与点$(-1, -1)$之间的距离，由几何关系可知该距离的最小值是1.
    \end{jieda}
    \item (2025$\cdot$嘉定区二模$\cdot$8)已知复数 ${z}_{1}$ 、 ${z}_{2}$ 满足 $\left| {z}_{1}\right|  = 1$ ， $\left| {z}_{2}\right|  = 2$ ， $\left| {{z}_{1} - {z}_{2}}\right|  = \sqrt{7}$ ，则 $\left| {{z}_{1} + {z}_{2}}\right|$ 的值为\blkda[2]{$\sqrt{3}$}.
    \begin{jieda}
        设 $\vv{OA}$ 对应的复数为 ${z}_{1},\vv{OB}$ 对应的复数为 ${z}_{2}$ ,则 $\vv{OA} + \vv{OB}$ 对应的复数为 ${z}_{1} + {z}_{2},\vv{OA} - \vv{OB}$ 对应的复数为 ${z}_{1} -  {z}_{2}$ ,因为 $\left| {z}_{1}\right|  = 1,\left| {z}_{2}\right|  = 2,\left| {{z}_{1} - {z}_{2}}\right|  = \sqrt{7}$ ,由平行四边形的性质可得: ${\left| {z}_{1} + {z}_{2}\right| }^{2} + {\left| {z}_{1} - {z}_{2}\right| }^{2} = 2\left( {{\left| {z}_{1}\right| }^{2} + {\left| {z}_{2}\right| }^{2}}\right)$ ,所以 $\left| {{z}_{1} + {z}_{2}}\right|  = \sqrt{2{\left| {z}_{1}\right| }^{2} + 2{\left| {z}_{2}\right| }^{2} - {\left| {z}_{1} - {z}_{2}\right| }^{2}} = \sqrt{2 + 8 - 7} =  \sqrt{3}$
    \end{jieda}
\end{enumerate}

\subsection*{三角换元与复数乘法的几何意义}
\begin{enumerate}
    \item 计算：$2\left(cos\dfrac{\pi}{9}+isin\dfrac{\pi}{9}\right) \cdot 3\left(cos\dfrac{\pi}{18} + i sin\dfrac{\pi}{18}\right)$
    \begin{jieda}
        \begin{align*}
            & 2\left(cos\dfrac{\pi}{9}+isin\dfrac{\pi}{9}\right) \cdot 3\left(cos\dfrac{\pi}{18} + i sin\dfrac{\pi}{18}\right) \\
            = & 6 \left(cos\dfrac{\pi}{6} + i sin\dfrac{\pi}{6}\right) \\
            = & 3\sqrt{3} + 3i
        \end{align*}
    \end{jieda}
    \item (2025$\cdot$黄浦区二模$\cdot$5)$i$为虚数单位，若复数 $z$ 满足 $z - \overline{z} = 2{i}$ 且 $\overline{z} = {i}z$ ，则 ${Re}z =$ \blkda[2]{$-1$}.
\end{enumerate}
\subsection*{乘方和开方}
\begin{enumerate}
    \item 计算：$(1+i)^{10}-(1-i)^{10}$
    \begin{jieda}
        设$z=(1+i) = \sqrt{2}\left(cos\dfrac{\pi}{4}+isin\dfrac{\pi}{4}\right)$，则$z^{10}=32\left(cos\dfrac{5\pi}{2}+isin\dfrac{5\pi}{2}\right) = 32i$ \\
        $(1-i)^{10} = \overline{(z^{10})} = (\overline{z})^{10} = -32i$ \\
        故目标值为$64i$
    \end{jieda}
    \item 已知$|z| = 1$，且$z^2+2z+\dfrac{1}{z}$是负实数，则复数$z = $\blkda[3]{$-\dfrac{1}{2} \pm \dfrac{\sqrt{3}}{2} i$或$-1$}
    \begin{jieda}
        $\dfrac{1}{z} = \dfrac{\overline{z}}{|z|^2}$，因为$|z| = 1$，所以$\dfrac{1}{z} = \overline{z}$，设$z = cos\alpha+i sin\alpha$，故$z^2+2z+\dfrac{1}{z} = cos2\alpha + i sin2\alpha + 2 cos\alpha + i(2sin\alpha) + cos\alpha - i(sin\alpha) = (cos2\alpha + 3cos\alpha) + i(sin\alpha + sin2\alpha)$ \\
        因为$z^2+2z+\dfrac{1}{z}$是负实数，所以$\left \{ \begin{array}{l}
            cos2\alpha + 3cos\alpha < 0   \\
            sin\alpha + sin2\alpha = 0  
        \end{array}\right.$，下式解得$sin\alpha = 0$或$cos\alpha = -\dfrac{1}{2}$，即$z = \pm1$或$-\dfrac{1}{2} \pm \dfrac{\sqrt{3}}{2} i$，代入上式检验可知$z = -\dfrac{1}{2} \pm \dfrac{\sqrt{3}}{2} i$或$-1$
    \end{jieda}
    \item 若$\omega = \dfrac{1+i}{2-2i}$是复数$z$的五次方根，求$z$的其余的五次方根.
    \begin{jieda}
        $\omega = \dfrac{1}{2} \cdot \dfrac{1+i}{1-i} = \dfrac{i}{2}$，因此$z = \dfrac{i}{32} = \dfrac{1}{32}\left(cos\dfrac{\pi}{2} + i sin\dfrac{\pi}{2}\right)$，故$z$的五次方根为$\dfrac{1}{2}\left(cos\left(\dfrac{\pi}{10}+\dfrac{2k\pi}{5}\right) + i sin\left(\dfrac{\pi}{10} + \dfrac{2k\pi}{5}\right)\right), k \in \{0,1,2,3,4\}$，其中$k = 1$对应$\omega$，因此其余的五次方根为$\dfrac{1}{2}\left(cos\left(\dfrac{\pi}{10}+\dfrac{2k\pi}{5}\right) + i sin\left(\dfrac{\pi}{10} + \dfrac{2k\pi}{5}\right)\right), k \in \{0,2,3,4\}$
    \end{jieda}
\end{enumerate}
\subsection*{乘积的模、商的模}
\begin{enumerate}
    \item (2024$\cdot$嘉定区二模$\cdot$5改)已知 ${i}$ 是虚数单位，$\left| \dfrac{a + {3{i}}}{{\left( 1 - 2{i}\right) }^{2}}\right|  = 1$，则实数$a = $ \blkda[2]{$\pm 4$}.
    \item (2025$\cdot$青浦区二模$\cdot$9)已知复数 $z$ 、 $w$ 满足 $\left| z\right|  = 2,z = \left( {1 + {i}}\right) w$ (i 是虚数单位)，则 $\left| {w + 2}\right|$ 的最大值是\blkda[2]{$2 + \sqrt{2}$}.
    \begin{jieda}
        已知 $\left| z\right|  = 2,z = \left( {1 + {i}}\right) w$ ,则 $\left| z\right|  = \left| {\left( {1 + {i}}\right) w}\right|  =  \mid  1 +  {i}\left| \right| w\left| \text{ . 因为 }\right| 1 + {i} \mid   = \sqrt{{1}^{2} + {1}^{2}} = \sqrt{2}$ ,所以 $\left| w\right|  = \dfrac{\left| z\right| }{\left| 1 + {i}\right| } = \dfrac{2}{\sqrt{2}} =  \sqrt{2},\left| {w + 2}\right|$ 表示复数 $w$ 所对应的点到 $\left( {-2,0}\right)$ 的距离, $\left| w\right|  =  \sqrt{2}$ 说明 $w$ 对应的点在以原点为圆心, $\sqrt{2}$ 为半径的圆上,所以 $\left| {w + 2}\right|$ 的最大值为圆心 $\left( {0,0}\right)$ 到点 $\left( {-2,0}\right)$ 的距离加上半径, 即 $2 + \sqrt{2}$.
    \end{jieda}
    \item (2024$\cdot$闵行区一模$\cdot$15)已知复数 ${z}_{1}$ 、 ${z}_{2}$ 在复平面内对应的点分别为 $P$ 、 $Q$ ， $\left| {OP}\right|  = 5$ ( $O$ 为坐标原点),且 ${z}_{1}^{2} - {z}_{1}{z}_{2} \cdot  \sin \theta  + {z}_{2}^{2} = 0$ ,则对任意 $\theta  \in  \mathbf{R}$ ,下列选项中为定值的是 \dotfill{(\kh{A})}
    \xx{$\left| {OQ}\right|$}{ $\left| {PQ}\right|$}{$\bigtriangleup {OPQ}$ 的周长}{$\bigtriangleup {OPQ}$ 的面积}
    \begin{jieda}
        $sin\theta = \dfrac{z_1}{z_2} + \dfrac{z_2}{z_1}$，因为$sin\theta \in \mathbf{R}$，所以$\dfrac{z_1}{z_2} = \overline{\left(\dfrac{z_2}{z_1}\right)}$，故$\left|\dfrac{z_1}{z_2}\right| = 1$，故$|z_2| = |z_1| = 5$.
    \end{jieda}
\end{enumerate}
\subsection*{乘积的共轭、商的共轭}
\begin{enumerate}
    \item (2025$\cdot$上海秋季卷$\cdot$10)已知复数 $z$ 满足 ${z}^{2} = {\left( \overline{z}\right) }^{2},\left| z\right|  \leq  1$ ，则 $\left| {z - 2 - {3i}}\right|$ 的最小值是\blkda[2]{$2\sqrt{2}$}.
    \begin{jieda}
        ${z}^{2} = {\left( \overline{z}\right) }^{2} = \overline{z^2}$，故$z^2 \in \mathbf{R}$，即$z$是实数或纯虚数.
        
        由复数的几何意义知 $z$ 在复平面内对应的点 $Z\left( {a,b}\right)$ 在单位圆内部(含边界)的坐标轴上运动， 如图所示即线段 ${AB},{CD}$ 上运动,设 $E\left( {2,3}\right)$ ,则 $\left| {z - 2 - 3{i}}\right|  = \left| {ZE}\right|$ ,由图像可知 $\left| {BE}\right|  = \sqrt{10} > \left| {CE}\right|  = 2\sqrt{2}$ ,所以 ${\left| ZE\right| }_{\min } = 2\sqrt{2}$.
        \begin{center}
            \includegraphics[width=0.2\textwidth]{bodyxb/src/104_1340_1138_236_243_0.jpg}
        \end{center}
    \end{jieda}

    \item (2023$\cdot$上海春季卷$\cdot$11)设 ${z}_{1},{z}_{2} \in  \mathbf{C}$ 且 ${z}_{1} = {i}\overline{{z}_{2}}$ ，满足 $\left| {{z}_{1} - 1}\right|  = 1$ ，则 $\left| {{z}_{1} - {z}_{2}}\right|$ 的取值范围为\blkda[2]{$\left\lbrack  {0,2 + \sqrt{2}}\right\rbrack$}.
    \begin{jieda}
        满足 $\left| {{z}_{1} - 1}\right|  = 1,{z}_{1}$ 以 $G(1,0)$ 为圆心,半径为 1 的圆
        
        ${z}_{1} = {i}\overline{{z}_{2}} \Longleftrightarrow z_2 = \overline{\left(\dfrac{z_1}{i}\right)}$，$z_2$是$z_1$逆时针旋转$\dfrac{\pi}{2}$再关于实轴对称，故${z}_{2}$ 以 $M\left( {0,1}\right)$ 为圆心,半径为 1 的圆,所以 $\left| {{z}_{1} - {z}_{2}}\right|$ 的取值范围 $\lbrack 0$ , $2 + \sqrt{2}\rbrack$ .

        \begin{center}
        \includegraphics[width=0.2\textwidth]{bodyxb/src/104_1343_1638_236_236_0.jpg}
        \end{center}
    \end{jieda}
\end{enumerate}
\subsection*{实系数一元二次方程}
\begin{enumerate}
    \item 若$-1+\sqrt{3}i$是实系数方程$x^2+ax+b=0$的根，则实数$a = $\blkda{2}，$b = $\blkda{4}.
    \begin{jieda}
        $-1+\sqrt{3}i$是方程一根，则$-1-\sqrt{3}i$是另一根，故根据韦达定理有$a = -(x_1 + x_2) = 2, b = x_1 \cdot x_2 = 4$
    \end{jieda}
    \item (2024$\cdot$黄浦区二模$\cdot$8)若实系数一元二次方程 ${x}^{2} + {ax} + b = 0$ 有一个虚数根的模为 4，则 $a$ 的取值范围是\blkda[2]{$\left( {-8,8}\right)$}.
    \begin{jieda}
        $16 = \left|z\right|^2 = z_1z_2 = b$，$\Delta = a^2 - 4b = a^2 - 64 < 0$，故$a \in (-8, 8)$
    \end{jieda}
    \item (2025$\cdot$虹口区二模$\cdot$8)已知 $z$ 是实系数一元二次方程 ${x}^{2} + {px} + q = 0$ 的一个虚根，且 $\left| {z - 1}\right|  = 2$ ， 若 $z$ 在复平面上所对应的点在抛物线 ${y}^{2} = {4x}$ 上,则 $p =$ \blkda[2]{$-2$}.
    \begin{jieda}
        $z$对应抛物线$y^2 = 4x$和圆$(x-1)^2+y^2 = 4$的交点，联立解得$x = 1$，即$Rez = 1$
        
        $p = -2Rez = -2$.
    \end{jieda}
\end{enumerate}